% GENERATED FILE. Edit canonical lessons, then run npm run generate:books.
% canonical-source-hash: fnv1a64:f4132f364fce1234
% canonical-lessons: JA-C09-sumimasen, JA-W09-mi, JA-W09-me, JA-C09-yoku, JA-W09-ku, JA-C09-wakarimasen, JA-W09-do, JA-C09-mou, JA-W09-mo, JA-C09-ichido, JA-C09-onegaishimasu, JA-C09-mou-ichido-onegaishimasu

\chapter{One More Time, Please}
\label{ch:ja-repeat-please}
\hlchaptermodality{9}

\begin{chapteropening}
\textbf{By the end of this chapter:} \emph{I can say that I do not understand, politely ask for one repetition, read and write the earned block \ja{もういちど}, and write five new hiragana forms without a chart.}

A spoken repair assembled from independently learned pieces, with decoding and dictation limited to the signs the learner has actually earned.
\end{chapteropening}

\section[sumimasen]{\ja{すみません} — open the repair gently}

\label{lesson:JA-C09-sumimasen}



\begin{quote}
Hear the farewell in your mind, say it, read \textbf{\ja{さようなら}}, and write it once. Now imagine that the other person has spoken and you missed it. Do not guess; open a repair.
\end{quote}

\subsection*{What to know first}
\begin{quote}
\textbf{\ja{すみません。}} — \emph{su-mi-ma-se-n} — excuse me; I am sorry
\end{quote}

Five morae, each given room. This polite expression can apologise, attract attention, or soften the request that follows. Here it means: “Excuse me — I need help with what was just said.”

Its shape is useful, not a literal translation burden. It is related to \emph{sumu}, “to be settled or finished,” in a polite negative form: the small social debt is not yet settled. Keep the practical job first.

\subsection*{Your turn}
\begin{itemize}
\raggedright
  \item \emph{Hear:} \emph{sumimasen} $\to$ point to \textbf{repair}, not \textbf{farewell}
  \item \emph{Say it:} \emph{su | mi | ma | se | n}{]} {[}REPEAT x2
  \item \emph{Choose:} you missed a word $\to$ \textbf{\ja{すみません}}
\end{itemize}

The Japanese spelling is exposure for now. The next two lessons isolate its two new signs before either becomes load-bearing.

\begin{tcolorbox}[breakable,colback=teal!4,colframe=teal!35!black,title={Before you move on}]
Say it once, then stop. The request itself comes later.

Source: \href{https://www.irodori.jpf.go.jp/assets/data/starter/pdf/X\_L02.pdf}{Japan Foundation Irodori Starter lesson 2}, pp. 1–3.
\end{tcolorbox}
\section[mi]{\ja{み} — one looped bridge}

\label{lesson:JA-W09-mi}



\begin{quote}
Say \emph{sumimasen} once, then write known \textbf{\ja{ま}} and \textbf{\ja{な}}. The new sign shares the \emph{m} opening but carries a different vowel.
\end{quote}

\subsection*{Script — the shape}
\begin{quote}
\textbf{\ja{み}} — \emph{mi}
\end{quote}

Write it in two strokes. First make the bent left body, ending in a small turn. Then begin at the upper right, sweep down through the middle, and finish by lifting to the right. Say \emph{mi} only when the second stroke ends.

\subsection*{Your turn}
1. Trace \textbf{\ja{み}} once with a finger. 2. Trace it once with a pencil. 3. Point to \textbf{\ja{み}} inside \textbf{\ja{すみません}}; do not decode the other signs yet.

\begin{tcolorbox}[breakable,colback=teal!4,colframe=teal!35!black,title={Before you move on}]
Cover the model and air-write the two strokes once.

Source: \href{https://www.unicode.org/charts/PDF/U3040.pdf}{Unicode Hiragana chart}; stroke order per \emph{Hitsujun Shido no Tebiki} (Japanese Ministry of Education, 1958).
\end{tcolorbox}
\section[me]{\ja{め} — cross, loop, release}

\label{lesson:JA-W09-me}



\begin{quote}
Find \textbf{\ja{み}} in \textbf{\ja{すみません}}, then say the five-mora repair opener.
\end{quote}

\subsection*{Script — the shape}
\begin{quote}
\textbf{\ja{め}} — \emph{me}
\end{quote}

Two strokes. The short first stroke leans down to the right. The longer second stroke crosses it, curves into an open loop, and releases at the lower right. Keep that final opening visible; do not close it into a circle.

\subsection*{Your turn}
\begin{itemize}
\raggedright
  \item \emph{Copy:} \textbf{\ja{め}} once with the model visible
  \item \emph{Read it:} point to \textbf{\ja{み}}, then \textbf{\ja{め}} and say \emph{mi, me}
  \item \emph{Find:} circle \textbf{\ja{め}} inside \textbf{\ja{すみません}}
\end{itemize}

\begin{tcolorbox}[breakable,colback=teal!4,colframe=teal!35!black,title={Before you move on}]
Write \textbf{\ja{め}} once more while saying \emph{me}.

Source: \href{https://www.unicode.org/charts/PDF/U3040.pdf}{Unicode Hiragana chart}; stroke order per \emph{Hitsujun Shido no Tebiki} (Japanese Ministry of Education, 1958).
\end{tcolorbox}
\section[yoku]{\ja{よく} — how fully?}

\label{lesson:JA-C09-yoku}



\begin{quote}
Say \emph{sumimasen}. Point to \textbf{\ja{み}} and \textbf{\ja{め}} inside it.
\end{quote}

\subsection*{What to know first}
\begin{quote}
\textbf{\ja{よく}} — \emph{yo-ku} — well; fully
\end{quote}

Two morae. \emph{Yoku} is the adverbial partner of \emph{yoi}, “good.” In the coming line, it says that understanding did not happen \textbf{well} or \textbf{fully}. It does not itself mean “understand.”

You already write \textbf{\ja{よ}}. Treat \textbf{\ja{く}} as visible support for now; its own tiny writing lesson follows.

\subsection*{Your turn}
\begin{itemize}
\raggedright
  \item \emph{Hear:} \emph{yoku} $\to$ choose \textbf{well/fully}
  \item \emph{Say it:} \emph{yo | ku}{]} {[}REPEAT x2
  \item \emph{Point to:} known \textbf{\ja{よ}}, then supported \textbf{\ja{く}}
\end{itemize}

\begin{tcolorbox}[breakable,colback=teal!4,colframe=teal!35!black,title={Before you move on}]
Say \emph{yoku} once without adding the rest of the sentence.

Source: \href{https://www.irodori.jpf.go.jp/assets/data/starter/pdf/X\_L02.pdf}{Japan Foundation Irodori Starter lesson 2}, p. 1.
\end{tcolorbox}
\section[ku]{\ja{く} — one angled stroke}

\label{lesson:JA-W09-ku}



\begin{quote}
Write known \textbf{\ja{よ}} and say \emph{yoku}.
\end{quote}

\subsection*{Script — the shape}
\begin{quote}
\textbf{\ja{く}} — \emph{ku}
\end{quote}

One stroke. Begin high on the right, sweep down-left to the point, then turn and sweep down-right. Keep it narrow and open. The turn is part of the same stroke.

\subsection*{Your turn}
1. Look at \textbf{\ja{く}} for five seconds. 2. Hide it and write the one angled stroke. 3. Reveal, check, then join known \textbf{\ja{よ}} + \textbf{\ja{く}} $\to$ \textbf{\ja{よく}}.

\begin{tcolorbox}[breakable,colback=teal!4,colframe=teal!35!black,title={Before you move on}]
Read \textbf{\ja{よく}} once without romanization.

Source: \href{https://www.unicode.org/charts/PDF/U3040.pdf}{Unicode Hiragana chart}; stroke order per \emph{Hitsujun Shido no Tebiki} (Japanese Ministry of Education, 1958).
\end{tcolorbox}
\section[wakarimasen]{\ja{わかりません} — I do not understand}

\label{lesson:JA-C09-wakarimasen}



\begin{quote}
Read \textbf{\ja{よく}}, then open a repair with \emph{sumimasen}.
\end{quote}

\subsection*{What to know first}
\begin{quote}
\textbf{\ja{わかりません。}} — \emph{wa-ka-ri-ma-se-n} — I do not understand
\end{quote}

Six morae. The \emph{wakari-} part carries “understand”; polite \textbf{-masen} makes the predicate negative. Japanese often leaves out “I” when the situation already makes the experiencer obvious.

Combine the two known pieces only in sound for now:

\begin{quote}
\emph{yoku wakarimasen} — I do not understand well / fully
\end{quote}

\subsection*{Your turn}
\begin{itemize}
\raggedright
  \item \emph{Hear:} \emph{wakarimasen} $\to$ choose \textbf{does not understand}
  \item \emph{Say it:} \emph{wa | ka | ri | ma | se | n}
  \item \emph{Repair:} \emph{yoku wakarimasen}
\end{itemize}

\begin{tcolorbox}[breakable,colback=teal!4,colframe=teal!35!black,title={Before you move on}]
Do not add an invented pronoun. The short Japanese line is complete here.

Source: \href{https://www.irodori.jpf.go.jp/assets/data/starter/pdf/X\_L02.pdf}{Japan Foundation Irodori Starter lesson 2}, pp. 1–2.
\end{tcolorbox}
\section[do]{\ja{ど} — voice a sign you already own}

\label{lesson:JA-W09-do}



\begin{quote}
Write \textbf{\ja{と}}. Add the two-stroke dakuten beside it. Say \emph{wakarimasen} once as the reason a repetition may be needed.
\end{quote}

\subsection*{Script — carry the voicing mark}
\begin{quote}
\textbf{\ja{と}} \emph{to} + \textbf{\ja{゛}} $\to$ \textbf{\ja{ど}} \emph{do}
\end{quote}

The base shape does not change. Write \textbf{\ja{と}} first, then add the two short dakuten strokes at the upper right. The mark voices the opening consonant: \emph{t} becomes \emph{d} while the vowel stays \emph{o}.

\subsection*{Your turn}
1. Copy \textbf{\ja{と} $\to$ \ja{ど}} once. 2. Cover the model and write the voiced sign from the heard mora \emph{do}. 3. Read \textbf{\ja{と} / \ja{ど}} as \emph{to / do} without changing the vowel.

\begin{tcolorbox}[breakable,colback=teal!4,colframe=teal!35!black,title={Before you move on}]
Write \textbf{\ja{ど}} once more from memory.

Source: \href{https://www.unicode.org/charts/PDF/U3040.pdf}{Unicode Hiragana chart}; voicing pattern and stroke order per \emph{Hitsujun Shido no Tebiki} (Japanese Ministry of Education, 1958).
\end{tcolorbox}
\section[mō]{\ja{もう} — one more time begins here}

\label{lesson:JA-C09-mou}



\begin{quote}
Say \emph{sumimasen, yoku wakarimasen}. Write \textbf{\ja{め}} and \textbf{\ja{ど}} once.
\end{quote}

\subsection*{What to know first}
\begin{quote}
\textbf{\ja{もう}} — \emph{mō} — again; one more
\end{quote}

Two morae: \emph{mo | o}. The macron in \emph{mō} hides the second beat, but the Japanese spelling shows it with final \textbf{\ja{う}}. In other settings this adverb can mean “already” or “more.” A request for repetition selects “again.”

\subsection*{Your turn}
\begin{itemize}
\raggedright
  \item \emph{Hear:} \emph{mō} in a repair $\to$ choose \textbf{again}
  \item \emph{Say it:} \emph{mo | o} and tap twice
  \item \emph{Point to:} \textbf{\ja{も} | \ja{う}}; \textbf{\ja{う}} owns the held second beat
\end{itemize}

\begin{tcolorbox}[breakable,colback=teal!4,colframe=teal!35!black,title={Before you move on}]
Say it once more, keeping both beats.

Source: \href{https://www.irodori.jpf.go.jp/assets/data/starter/pdf/X\_L02.pdf}{Japan Foundation Irodori Starter lesson 2}, pp. 1–3.
\end{tcolorbox}
\section[mo]{\ja{も} — two bars, then the curve}

\label{lesson:JA-W09-mo}



\begin{quote}
Write known \textbf{\ja{う}} and \textbf{\ja{み}}, then say \emph{mo | o}.
\end{quote}

\subsection*{Script — the shape}
\begin{quote}
\textbf{\ja{も}} — \emph{mo}
\end{quote}

Three strokes. Draw the central curve first, from the upper left down and around. Then add the upper short horizontal and the lower short horizontal. Both bars cross the central stroke.

\subsection*{Your turn}
1. Copy \textbf{\ja{も}} once. 2. Hide it for ten seconds and write it again. 3. Add known \textbf{\ja{う}}: \textbf{\ja{もう}}. Read \emph{mo | o}.

\begin{tcolorbox}[breakable,colback=teal!4,colframe=teal!35!black,title={Before you move on}]
Cover the model and write \textbf{\ja{もう}} once.

Source: \href{https://www.unicode.org/charts/PDF/U3040.pdf}{Unicode Hiragana chart}; stroke order per \emph{Hitsujun Shido no Tebiki} (Japanese Ministry of Education, 1958).
\end{tcolorbox}
\section[ichido]{\ja{いちど} — one time, all known signs}

\label{lesson:JA-C09-ichido}



\begin{quote}
Write \textbf{\ja{い}}, \textbf{\ja{ち}}, and \textbf{\ja{と}}. Add the dakuten to \textbf{\ja{と}} and read the result \emph{do}. Write \textbf{\ja{め}} once as a review pulse.
\end{quote}

\subsection*{What to know first}
\begin{quote}
\textbf{\ja{いちど}} — \emph{i-chi-do} — one time; once
\end{quote}

The word is built transparently: \emph{ichi} is “one,” and \emph{do} counts one occurrence. All four visible marks are already yours: \textbf{\ja{い} + \ja{ち} + \ja{と} + \ja{゛}}.

Now join it to the previous word:

\begin{quote}
\textbf{\ja{もういちど}} — \emph{mō ichido} — one more time
\end{quote}

\subsection*{Your turn}
1. Hear \emph{i | chi | do} without looking. 2. Write \textbf{\ja{いちど}} from the three morae. 3. Read \textbf{\ja{もういちど}}, keeping the two beats in \textbf{\ja{もう}}.

\begin{tcolorbox}[breakable,colback=teal!4,colframe=teal!35!black,title={Before you move on}]
Circle the dakuten. It changes the final mora, not the first two.

Source: \href{https://www.irodori.jpf.go.jp/assets/data/starter/pdf/X\_L02.pdf}{Japan Foundation Irodori Starter lesson 2}, pp. 1–3.
\end{tcolorbox}
\section[onegai shimasu]{\ja{おねがいします} — please, as a polite request}

\label{lesson:JA-C09-onegaishimasu}



\begin{quote}
Read \textbf{\ja{いちど}}. Say \emph{sumimasen, yoku wakarimasen}. Write \textbf{\ja{く}}.
\end{quote}

\subsection*{What to know first}
\begin{quote}
\textbf{\ja{おねがいします。}} — \emph{o-ne-ga-i shi-ma-su} — please; I request it
\end{quote}

The pieces make the social direction visible: respectful \emph{o-}, \emph{negai} “request or wish,” and polite \emph{shimasu} “do.” Learn the whole as the courteous request ending. Do not translate each piece while speaking.

The common kanji spelling can wait. This chapter earns the sound and social use first, without switching writing systems mid-ramp.

\subsection*{Your turn}
\begin{itemize}
\raggedright
  \item \emph{Hear:} \emph{onegaishimasu} $\to$ choose \textbf{a polite request}
  \item \emph{Say it:} \emph{onegaishimasu} once at an even pace
  \item \emph{Join:} \emph{mō ichido, onegaishimasu}
\end{itemize}

\begin{tcolorbox}[breakable,colback=teal!4,colframe=teal!35!black,title={Before you move on}]
Keep the tone even and courteous; volume does not make the request more polite.

Source: \href{https://www.irodori.jpf.go.jp/assets/data/starter/pdf/X\_L02.pdf}{Japan Foundation Irodori Starter lesson 2}, pp. 1–3.
\end{tcolorbox}
\section[mō ichido, onegai shimasu]{\ja{もういちど、おねがいします} — ask for one more pass}

\label{lesson:JA-C09-mou-ichido-onegaishimasu}



\begin{quote}
Write \textbf{\ja{み}, \ja{め}, \ja{く}, \ja{ど}, \ja{も}}. Say \emph{sumimasen, yoku wakarimasen}. Then retrieve \emph{mō}, \emph{ichido}, and \emph{onegaishimasu} separately.
\end{quote}

\subsection*{What to know first}
\begin{quote}
\textbf{\ja{もういちど、おねがいします。}}
\end{quote}

>

\begin{quote}
\emph{mō ichido, onegai shimasu} — one more time, please
\end{quote}

Nothing new is hidden inside the line. \textbf{\ja{もう}} supplies “again,” \textbf{\ja{いちど}} supplies “one time,” and \textbf{\ja{おねがいします}} makes the request polite. The first two words are now readable; the final expression remains supported exposure until its own longer script runway is complete.

\subsection*{Your turn}
1. \textbf{Listen:} distinguish \emph{sumimasen} from \emph{mō ichido, onegaishimasu}. 2. \textbf{Speak:} say the opener, then the request, with a small pause. 3. \textbf{Read:} independently read \textbf{\ja{もういちど}}; use the printed support for the rest. 4. \textbf{Write:} from dictation, write only the independently earned block \textbf{\ja{もういちど}}.

\begin{tcolorbox}[breakable,colback=teal!4,colframe=teal!35!black,title={Before you move on}]
Stop after one accurate repair. Chapter 10 will ask for slower speech instead.

Source: \href{https://www.irodori.jpf.go.jp/assets/data/starter/pdf/X\_L02.pdf}{Japan Foundation Irodori Starter lesson 2}, pp. 1–3.
\end{tcolorbox}
